% q0057.tex — Surviving Losses
\question{
  id        = {0057},
  title     = {Surviving Losses},
  source    = {OBECON},
  year      = {2026},
  round     = {Seletiva IEO -- Phase C},
  language  = {English},
  topic     = {Micro},
  subtopic  = {Firm Theory / Shutdown Decision},
  type      = {objective},
  statement = {%
    Can a firm rationally operate at a loss in the short run? Under what
    condition should a loss-making firm shut down operations?
  },
  choices   = {%
    \item Not possible; operating at a loss is always worse than shutting
          down immediately.
    \item Not possible; a firm should shut down whenever total revenue falls
          below total costs.
    \item Possible if total revenue exceeds fixed costs; the firm should shut
          down when price falls below average total cost.
    \item Possible when fixed costs are sunk; the firm should shut down if
          price falls below average variable cost.
  },
  answer    = {D},
  solution  = {%
    In the short run, fixed costs are already sunk and irrelevant to the
    shutdown decision. The firm covers its avoidable (variable) costs as long
    as $P \geq AVC$, thereby reducing its loss compared to shutting down and
    bearing fixed costs alone. Shutdown occurs when $P < AVC$, because
    producing would add more to costs than to revenue. Long-run exit occurs
    when $P < ATC$, since all costs become variable in the long run. (D) is
    correct.
  },
}
